\hsection{Chapter 7}
\sol{703} $\log a+\log b=\log(ab)$ with $ab=(49+\sqrt{49^2-1})(49-\sqrt{49^2-1})=49^2-(49^2-1)=1$: the sum is $\log1=0$.

\sol{704} Put $r=e^{f(x)}=\dfrac{e^x+1}{e^x-1}$. Then $f(f(x))=\log\dfrac{r+1}{r-1}$ with $r+1=\dfrac{2e^x}{e^x-1}$ and $r-1=\dfrac2{e^x-1}$, so $\dfrac{r+1}{r-1}=e^x$ and $f(f(x))=x$. (Strictly, $f(x)$ is defined when $\frac{e^x+1}{e^x-1}>0$, i.e.\ $x>0$; then $f(x)>0$ as well, so $f\circ f$ makes sense there. For $x<0$ one reads $\log|\cdot|$.)

\sol{707} Both graphs pass through $(0,0)$: $\sin0=0=\arctan0$. The slopes there: $\cos0=1$ and $\frac1{1+0^2}=1$. Same point, same slope: the common tangent line is $y=x$.

\sol{708} $\dfrac{1+x^2}{1-x}=(1+x^2)(1+x+x^2+x^3+x^4+\dots)=1+x+2x^2+2x^3+2x^4+\dots$, so $P_4(x)=1+x+2x^2+2x^3+2x^4$.

\sol{709} $f=\arctan$ satisfies $f'=\frac1{1+x^2}$, i.e.\ $(1+x^2)f'=1$. As in 316, differentiate this relation: $2xf'+(1+x^2)f''=0$, $2f'+4xf''+(1+x^2)f'''=0$, $6f''+6xf'''+(1+x^2)f''''=0$, $12f'''+8xf''''+(1+x^2)f^{(5)}=0$. At $0$: $f'=1$, $f''=0$, $f'''=-2$, $f''''=0$, $f^{(5)}=24$. So $P_5(x)=x-\frac2{3!}x^3+\frac{24}{5!}x^5=x-\frac13x^3+\frac15x^5$. (Faster: $f'=1-x^2+x^4-\dots$, integrate term by term.)

\sol{711} (a) $x\xrightarrow{p}1+x\xrightarrow{p}2+x\xrightarrow{i}\frac1{2+x}\xrightarrow{t}\frac3{2+x}\xrightarrow{m}\frac{-3}{2+x}$: $m\circ t\circ i\circ p\circ p$. (b) $\dfrac{1-x}{2+x}=\dfrac{3-(2+x)}{2+x}=\dfrac3{2+x}-1$. From $g=\frac{-3}{2+x}$: $p(g(x))=1-\frac3{2+x}$ and $m(p(g(x)))=\frac3{2+x}-1$. So $\dfrac{1-x}{2+x}=m\circ p\circ m\circ t\circ i\circ p\circ p$.

\sol{712} With $\varphi=\arctan x$ in a right triangle with legs $1$ (adjacent) and $x$ (opposite), hypotenuse $\sqrt{1+x^2}$: $\sin(\arctan x)=\dfrac x{\sqrt{1+x^2}}$ (correct sign for $x<0$ too, since $\arctan x\in(-\frac\pi2,\frac\pi2)$ and $\sin$ has the sign of $x$ there).

\sol{714} $x=\cosh t=\frac12(e^t+e^{-t})$; with $u=e^t$: $2x=u+\frac1u$, $u^2-2xu+1=0$, $u=x\pm\sqrt{x^2-1}$. For $x\ge1$ both roots are positive (their product is $1$), so $t=\log\big(x\pm\sqrt{x^2-1}\big)$: two values, because $\cosh$ is not injective.

\sol{715} (a) If $f$ is even, then $y=f(x)$ has the solutions $x_0$ and $-x_0$; the inverse on $x\ge0$ returns $x_0$, the inverse on $x\le0$ returns $-x_0$. The two inverse graphs are therefore each other's image under $y\mapsto-y$: reflection in the $x$-axis. (b) Mirror images means $\log\big(x-\sqrt{x^2-1}\big)=-\log\big(x+\sqrt{x^2-1}\big)=\log\dfrac1{x+\sqrt{x^2-1}}$, i.e.\ $\dfrac1{x-\sqrt{x^2-1}}=x+\sqrt{x^2-1}$. True: $\big(x-\sqrt{x^2-1}\big)\big(x+\sqrt{x^2-1}\big)=x^2-(x^2-1)=1$.

\sol{716} (a) $\dfrac d{dx}\log\big(x+\sqrt{x^2-1}\big)=\dfrac{1+\frac x{\sqrt{x^2-1}}}{x+\sqrt{x^2-1}}=\dfrac1{\sqrt{x^2-1}}$ (the numerator is $\frac{\sqrt{x^2-1}+x}{\sqrt{x^2-1}}$). For the other branch: $\dfrac{1-\frac x{\sqrt{x^2-1}}}{x-\sqrt{x^2-1}}=-\dfrac1{\sqrt{x^2-1}}$. (b) Differentiate $\cosh(\operatorname{arcosh}x)=x$: $\sinh(\operatorname{arcosh}x)\cdot\operatorname{arcosh}'x=1$, and by Oefening 418 $\sinh t=\pm\sqrt{\cosh^2t-1}=\pm\sqrt{x^2-1}$ ($+$ for $t\ge0$, $-$ for $t\le0$), so $\operatorname{arcosh}'x=\pm\dfrac1{\sqrt{x^2-1}}$, matching (a).

\sol{717} $\tanh t=\dfrac{\sinh t}{\cosh t}=\dfrac{e^t-e^{-t}}{e^t+e^{-t}}$: odd, $\tanh0=0$, increasing, $\tanh t\to\pm1$ as $t\to\pm\infty$ (divide by $e^{t}$): an S-shaped graph between the asymptotes $y=\pm1$. Derivative (quotient rule and 418): $\tanh't=\dfrac{\cosh^2t-\sinh^2t}{\cosh^2t}=\dfrac1{\cosh^2t}=1-\tanh^2t$ (two forms, like $\tan'=\frac1{\cos^2}=1+\tan^2$). Inverse: $x=\tanh t=\dfrac{u-1/u}{u+1/u}=\dfrac{u^2-1}{u^2+1}$ with $u=e^t$, so $u^2=\dfrac{1+x}{1-x}$ and $t=\dfrac12\log\dfrac{1+x}{1-x}$ for $-1<x<1$; its derivative is $\dfrac12\Big(\dfrac1{1+x}+\dfrac1{1-x}\Big)=\dfrac1{1-x^2}$, consistent with $\operatorname{artanh}'x=\frac1{1-\tanh^2(\operatorname{artanh}x)}$.

\sol{721} For $x\ne1$: $\dfrac{1-x^n}{1-x}=1+x+x^2+\dots+x^{n-1}$ (Oefening 330a). Integrating term by term over $[0,1]$: $\sum_{k=0}^{n-1}\frac1{k+1}=1+\frac12+\dots+\frac1n$.

\sol{725} $\sqrt x$: domain $[0,\infty)$, through $(0,0)$ and $(1,1)$, increasing, concave, vertical tangent at $0$, $\to\infty$. $e^x$: domain $\R$, positive, $e^0=1$, increasing, convex, $\to0$ as $x\to-\infty$ (asymptote $y=0$), $\to\infty$ faster than any power. $\log x$: domain $(0,\infty)$, zero at $1$, $\to-\infty$ at $0^+$ (vertical asymptote), increasing, concave, $\to\infty$ slower than any power; mirror image of $e^x$ in $y=x$. $\sin x$: odd, period $2\pi$, zeros $k\pi$, extrema $\pm1$ at $\frac\pi2+k\pi$, inflection points at the zeros ($\sin''=-\sin$). $\cos x$: even, $\cos x=\sin(x+\frac\pi2)$. $\tan x$: odd, period $\pi$, domain $x\ne\frac\pi2+k\pi$ (vertical asymptotes), zeros $k\pi$, increasing on each branch ($\tan'=1+\tan^2>0$), inflection points at the zeros.

\sol{726} Each: domain; symmetry; intercepts; edges; $f'$; $f''$.\\
(a) $k(x)=\sqrt{9-x^2}$: domain $[-3,3]$, even, zeros $\pm3$, $k(0)=3$; $k'=-\frac x{\sqrt{9-x^2}}$: maximum $3$ at $0$, vertical tangents at $\pm3$; $k''=-\frac9{(9-x^2)^{3/2}}<0$: concave. The upper half of the circle of radius $3$.\\
(b) $h(x)=\frac x{\sqrt{1-x^2}}$: domain $(-1,1)$, odd, zero at $0$; $h\to\pm\infty$ as $x\to\pm1$ (vertical asymptotes); $h'=(1-x^2)^{-3/2}>0$, increasing, $h'(0)=1$; $h''=3x(1-x^2)^{-5/2}$: inflection at $0$, convex for $x>0$.\\
(c) $g(x)=e^{1-x^2}$: domain $\R$, even, no zeros, $g(0)=e$, $g\to0$ at $\pm\infty$ (asymptote $y=0$); $g'=-2xe^{1-x^2}$: maximum $e$ at $0$; $g''=(4x^2-2)e^{1-x^2}$: inflection points at $x=\pm\frac12\sqrt2$ (value $\sqrt e$); bell shape.\\
(e) $\varphi(x)=\log\frac{2x}{(x-4)^2}=\log2+\log x-2\log|x-4|$: domain $x>0$, $x\ne4$; zeros where $2x=(x-4)^2$: $x=2$ and $x=8$; $\varphi\to-\infty$ as $x\to0^+$ and as $x\to\infty$ ($\frac{2x}{(x-4)^2}\to0$), $\varphi\to+\infty$ as $x\to4$ (vertical asymptote); $\varphi'=\frac1x-\frac2{x-4}=-\frac{x+4}{x(x-4)}$: positive on $(0,4)$, negative on $(4,\infty)$, no critical points; $\varphi''=\frac{x^2+8x-16}{x^2(x-4)^2}$: zero at $x=4\sqrt2-4\approx1.66$, concave before it, convex after (on both sides of $4$); inflection point at $x=4\sqrt2-4$.\\
(f) $\psi(x)=(1+\frac1x)e^{-x}$: domain $x\ne0$; zero at $-1$; $\psi\to+\infty$ as $x\to0^+$, $-\infty$ as $x\to0^-$, $0^+$ as $x\to\infty$, $+\infty$ as $x\to-\infty$; $\psi'=-\frac{x^2+x+1}{x^2}e^{-x}<0$ everywhere: decreasing on each branch; $\psi''=\frac{(x+1)(x^2+2)}{x^3}e^{-x}$: convex for $x<-1$ and for $x>0$, concave on $(-1,0)$; inflection point $(-1,0)$.\\
(g) $z(x)=e^{-1/x}$: domain $x\ne0$, positive; $z\to0$ as $x\to0^+$, $\to+\infty$ as $x\to0^-$, $\to1$ as $x\to\pm\infty$ (asymptote $y=1$, approached from below on the right and from above on the left); $z'=\frac1{x^2}e^{-1/x}>0$, increasing on each branch; $z''=\frac{1-2x}{x^4}e^{-1/x}$: convex for $x<\frac12$ ($x\ne0$), concave for $x>\frac12$, inflection at $(\frac12,e^{-2})$.\\
(h) $L(x)=\log(x-\frac1x)$: domain where $\frac{x^2-1}x>0$: $(-1,0)\cup(1,\infty)$; zeros where $x^2-x-1=0$: $x=\frac{1\pm\sqrt5}2$ (one in each interval); $L\to-\infty$ at $-1^+$ and $1^+$, $+\infty$ at $0^-$ and $+\infty$; $L'=\frac{x^2+1}{x(x^2-1)}>0$ on the domain: increasing; $L''=-\frac{x^4+4x^2-1}{x^2(x^2-1)^2}$: zero at $x=\pm\sqrt{\sqrt5-2}\approx\pm0.486$; on $(-1,-0.486)$ convex, on $(-0.486,0)$ concave (inflection at $-0.486$); on $(1,\infty)$ concave.\\
(i) $e^{1-\frac1{1+x}}$: domain $x\ne-1$, positive; $\to0$ as $x\to-1^+$, $\to+\infty$ as $x\to-1^-$, $\to e$ as $x\to\pm\infty$ (asymptote $y=e$); derivative $\frac1{(1+x)^2}e^{1-\frac1{1+x}}>0$: increasing on each branch; second derivative $-\frac{2x+1}{(1+x)^4}e^{1-\frac1{1+x}}$: convex for $x<-\frac12$ ($x\ne-1$), concave for $x>-\frac12$, inflection at $x=-\frac12$ (value $e^{-1}$).\\
(j) $\arcsin(1-\sqrt x)$: domain $0\le x\le4$ ($-1\le1-\sqrt x\le1$); values from $\frac\pi2$ at $0$ through $0$ at $x=1$ to $-\frac\pi2$ at $4$; derivative $-\frac1{2\sqrt x\sqrt{2\sqrt x-x}}<0$: decreasing, vertical tangents at $0$ and $4$; second derivative changes sign at $\sqrt x=\frac32$, i.e.\ $x=\frac94$: convex on $(0,\frac94)$, concave on $(\frac94,4)$.

\sol{730} Let $h(x)=2(\sqrt x-1)-\log x$ on $(0,\infty)$. $h(1)=0$ and $h'(x)=\frac1{\sqrt x}-\frac1x=\frac{\sqrt x-1}x$, negative on $(0,1)$, positive on $(1,\infty)$. So $h$ has its minimum $0$ at $x=1$ and $h\ge0$: $\log x\le2(\sqrt x-1)$. (This is sharper than 729 for $x$ near $1$, since $2(\sqrt x-1)\le x-1$ by $(\sqrt x-1)^2\ge0$.)
